\documentclass[aspectratio=169]{beamer}

% Shared Beamer settings for Elements of AIML lectures (UPES).
% Keep this file free of \documentclass to allow reuse via \input.

% Allow lecture sources to use shared theme/assets without copying them into each lecture folder.
% NOTE: These paths are relative to `Unit-xx/Lecture-yy/latex/` (and `Lab/Experiment-xx/latex/`).
\makeatletter
\def\input@path{{../../../_shared/}{../../../_shared/UPES-Beamer-Theme/}}
\makeatother

\usetheme{UPES}
\setbeamertemplate{navigation symbols}{}

\usepackage{amsmath}
\usepackage{amssymb}
\usepackage{booktabs}
\usepackage{graphicx}
\usepackage{hyperref}
\usepackage{xcolor}
\usepackage{listings}

% Code listings (general-purpose).
\lstset{
  basicstyle=\ttfamily\small,
  keywordstyle=\color{blue},
  commentstyle=\color{gray},
  stringstyle=\color{teal},
  showstringspaces=false,
  columns=fullflexible,
  breaklines=true
}

\usepackage{tikz}
\usetikzlibrary{arrows.meta,positioning,shapes.geometric,calc}

% Per-slide source line (references shown on the slide itself).
\newcommand{\slidesources}[1]{%
  \vfill
  {\tiny\color{gray}\raggedright\sloppy\parbox{\linewidth}{Sources: #1}\par}%
}

% Unit-wise source strings for this subject (used by slides/notes).
% Path is relative to the lecture `latex/` folder (the compilation CWD).
% Source strings for Elements of AIML (used by slides + notes).

\newcommand{\AIMLRepoURL}{https://github.com/tali7c/Elements-of-AIML}

\newcommand{\AIMLLabSources}{%
Provost and Fawcett, \textit{Data Science for Business}.;%
McKinney, \textit{Python for Data Analysis}.;%
WEKA Documentation (Waikato Environment for Knowledge Analysis), accessed 2026-02-28.;%
Matplotlib and Seaborn documentation, accessed 2026-02-28.%
}

\newcommand{\AIMLUnitISources}{%
Rich and Knight, \textit{Artificial Intelligence}, McGraw Hill, 2017.;%
Russell and Norvig, \textit{Artificial Intelligence: A Modern Approach} (4e), Ch. 1--2 (Introduction, Intelligent Agents).;%
Poole and Mackworth, \textit{Artificial Intelligence: Foundations of Computational Agents} (2e), selected sections.%
}

\newcommand{\AIMLUnitIISources}{%
Russell and Norvig, \textit{Artificial Intelligence: A Modern Approach} (4e), Ch. 7--9 (Logical Agents, First-Order Logic, Inference).;%
Huth and Ryan, \textit{Logic in Computer Science} (2e), selected sections on propositional and first-order logic.;%
Genesereth and Nilsson, \textit{Logical Foundations of Artificial Intelligence}, selected chapters.%
}

\newcommand{\AIMLUnitIIISources}{%
Mueller and Massaron, \textit{Machine Learning for Dummies}, For Dummies, 2016.;%
Mitchell, \textit{Machine Learning}, selected chapters on learning basics and evaluation.;%
Scikit-learn documentation, accessed 2026-02-28.%
}

\newcommand{\AIMLUnitIVSources}{%
Mueller and Massaron, \textit{Machine Learning for Dummies}, For Dummies, 2016.;%
James, Witten, Hastie, and Tibshirani, \textit{An Introduction to Statistical Learning} (2e), selected chapters.;%
Scikit-learn documentation, accessed 2026-02-28.%
}

\newcommand{\AIMLUnitVSources}{%
NIST, \textit{AI Risk Management Framework (AI RMF 1.0)}, 2023.;%
Barocas, Hardt, and Narayanan, \textit{Fairness and Machine Learning} (online book), selected sections.;%
Knaflic, \textit{Storytelling with Data}, selected chapters on communicating results.%
}



\title[Elements of AIML]{Elements of AIML}
\subtitle{Miscellaneous -- Lecture 02: Eigenvalues and Eigenvectors from a Data Sample}
\author{Tofik Ali}
\institute{School of Computer Science, UPES Dehradun}
\date{\today}

\graphicspath{{../images/}{../../../_shared/}{../../../_shared/UPES-Beamer-Theme/}}

\begin{document}

\UPESCoverFrame
\setcounter{framenumber}{0}

\begin{frame}
  \titlepage
  \vspace{0.3em}
  \begin{center}
    \footnotesize Supplementary derivation lecture: \textbf{from raw sample data to covariance eigenvectors}\\[0.25em]
    \footnotesize Repository: \texttt{\AIMLRepoURL}
  \end{center}
  \slidesources{\AIMLMiscSources}
\end{frame}

\begin{frame}{Agenda}
  \tableofcontents
\end{frame}

\section{Context}

\begin{frame}{Learning Outcomes}
  By the end of this lecture, you should be able to:
  \begin{itemize}[<+->]
    \item start from a small raw dataset and compute the mean vector,
    \item center the data and build the covariance matrix,
    \item solve for eigenvalues and eigenvectors by hand,
    \item interpret the largest eigenvector as the first PCA direction.
  \end{itemize}
\end{frame}

\begin{frame}{Why This Matters}
  \begin{itemize}[<+->]
    \item In \textbf{PCA}, eigenvectors define the principal directions.
    \item In \textbf{LDA}, a related eigen-problem appears using scatter matrices.
    \item This lecture shows the mechanics on a tiny dataset where hand calculation is still realistic.
  \end{itemize}
\end{frame}

\section{Raw Sample}

\begin{frame}{Raw Sample Data}
  Consider three two-dimensional samples:
  \[
    x_1 =
    \begin{bmatrix}
      2\\1
    \end{bmatrix},
    \quad
    x_2 =
    \begin{bmatrix}
      1\\2
    \end{bmatrix},
    \quad
    x_3 =
    \begin{bmatrix}
      0\\0
    \end{bmatrix}
  \]
  \vspace{0.3em}
  Goal: compute eigenvalues and eigenvectors \textbf{from this sample}, not from an already given covariance matrix.
\end{frame}

\section{Centering}

\begin{frame}{Step 1: Compute the Mean}
  Mean vector:
  \[
    \mu =
    \frac{1}{3}
    \left(
    \begin{bmatrix}
      2\\1
    \end{bmatrix}
    +
    \begin{bmatrix}
      1\\2
    \end{bmatrix}
    +
    \begin{bmatrix}
      0\\0
    \end{bmatrix}
    \right)
    =
    \begin{bmatrix}
      1\\1
    \end{bmatrix}
  \]
\end{frame}

\begin{frame}{Step 2: Center the Data}
  Subtract the mean from each sample:
  \[
    x_1 - \mu =
    \begin{bmatrix}
      1\\0
    \end{bmatrix},
    \quad
    x_2 - \mu =
    \begin{bmatrix}
      0\\1
    \end{bmatrix},
    \quad
    x_3 - \mu =
    \begin{bmatrix}
      -1\\-1
    \end{bmatrix}
  \]
  So the centered data matrix is:
  \[
    X_c =
    \begin{bmatrix}
      1 & 0\\
      0 & 1\\
      -1 & -1
    \end{bmatrix}
  \]
\end{frame}

\section{Covariance}

\begin{frame}{Step 3: Build the Covariance Matrix}
  \[
    C = \frac{1}{n-1} X_c^\top X_c
  \]
  Here $n=3$, so:
  \[
    X_c^\top X_c =
    \begin{bmatrix}
      1 & 0 & -1\\
      0 & 1 & -1
    \end{bmatrix}
    \begin{bmatrix}
      1 & 0\\
      0 & 1\\
      -1 & -1
    \end{bmatrix}
    =
    \begin{bmatrix}
      2 & 1\\
      1 & 2
    \end{bmatrix}
  \]
  Therefore:
  \[
    C =
    \begin{bmatrix}
      1 & 0.5\\
      0.5 & 1
    \end{bmatrix}
  \]
\end{frame}

\section{Eigenvalues}

\begin{frame}{Step 4: Solve the Characteristic Equation}
  We solve:
  \[
    \det(C - \lambda I)=0
  \]
  \[
    \det
    \begin{bmatrix}
      1-\lambda & 0.5\\
      0.5 & 1-\lambda
    \end{bmatrix}
    = 0
  \]
  \[
    (1-\lambda)^2 - 0.25 = 0
  \]
  \[
    \lambda^2 - 2\lambda + 0.75 = 0
  \]
  So the eigenvalues are:
  \[
    \lambda_1 = 1.5,
    \qquad
    \lambda_2 = 0.5
  \]
\end{frame}

\section{Eigenvectors}

\begin{frame}{Step 5: Solve for Eigenvectors}
  For $\lambda_1 = 1.5$:
  \[
    C - 1.5I =
    \begin{bmatrix}
      -0.5 & 0.5\\
      0.5 & -0.5
    \end{bmatrix}
  \]
  Solving $(C-1.5I)v=0$ gives $y=x$, so:
  \[
    v_1 \propto
    \begin{bmatrix}
      1\\1
    \end{bmatrix}
  \]
  For $\lambda_2 = 0.5$:
  \[
    v_2 \propto
    \begin{bmatrix}
      1\\-1
    \end{bmatrix}
  \]
\end{frame}

\begin{frame}{Step 6: Normalize the Eigenvectors}
  Normalized versions:
  \[
    \hat v_1 =
    \frac{1}{\sqrt{2}}
    \begin{bmatrix}
      1\\1
    \end{bmatrix},
    \qquad
    \hat v_2 =
    \frac{1}{\sqrt{2}}
    \begin{bmatrix}
      1\\-1
    \end{bmatrix}
  \]
  \vspace{0.3em}
  Since $\lambda_1 > \lambda_2$, $\hat v_1$ is the first principal direction for PCA.
\end{frame}

\section{Interpretation}

\begin{frame}{Interpretation for PCA}
  \begin{itemize}[<+->]
    \item The first principal direction is the diagonal line $x_2 = x_1$.
    \item The first eigenvalue is larger, so that direction captures more variance.
    \item The second direction is perpendicular and captures the remaining smaller spread.
  \end{itemize}
\end{frame}

\begin{frame}{Common Mistakes}
  \begin{itemize}[<+->]
    \item Using raw data directly without centering it.
    \item Taking eigenvalues from the data matrix instead of the covariance matrix.
    \item Forgetting to normalize the eigenvectors.
    \item Forgetting that the largest eigenvalue corresponds to the first PCA component.
  \end{itemize}
\end{frame}

\section{Summary}

\begin{frame}{Summary}
  \begin{itemize}
    \item Start from the raw sample, compute the mean, and center the data.
    \item Build the covariance matrix before solving the eigen-problem.
    \item Solve $\det(C-\lambda I)=0$, then solve $(C-\lambda I)v=0$.
    \item In PCA, the eigenvector of the largest eigenvalue becomes the first principal component.
  \end{itemize}
  \slidesources{\AIMLMiscSources}
\end{frame}

\begin{frame}{Exit Questions}
  \begin{enumerate}
    \item Why do we center the data before computing covariance?
    \item What matrix do we use for the eigenvalue problem in PCA?
    \item Which eigenvector becomes PC1?
    \item Why must eigenvectors usually be normalized?
  \end{enumerate}
\end{frame}

\end{document}
